\section{Literature Review}
%Previous studies summary}
%What pertinent literature needs to be reviewed to orient the reader and provide sufficient background information for them to understand your study?.
ASTM-F standards and mechanical risks
\cite{Schenck2000} safety
\cite{Ghadimi2025} Contraindications

both
\cite{stoianovici2024}, 

force
\cite{Yin2024473}

torque
\cite{Seo2021767}


%What are the main arguments or findings from the sources most relevant to your argument, especially those you intend to challenge, modify, or expand upon?.

Do I need to add a case when a metal object was attracted violently towards the bore? Just to make an impact... 

%Have there been previous effective reviews of your subject, and if so, what is the starting point (e.g., date of previous review) for your current review?.
Previous studies made in our institution seek to illustrate the interaction of leads with the \gls{rf} field in terms of energy absorbed and heat produced. It showed that lead length, type, and orientation significantly influence heating behavior.This study found that short leads ($<13$ cm) posed no significant thermal hazard during MRI at a specific absorption rate (SAR) of 2 W/kg and is throughtfully described as Master Thesis in 2021 \cite{haddixProposal,astmF2182,aboyewa2021}.

Along the same line of thought, we aim to cover the mechanical aspect of lead-\gls{mri} interaction covering (1) Characterization of the static magnetic field, (2) quantification of translational forces near the bore entrance, (3) Torque evaluation from bore entrance to isocenter.

While some preliminary work has been conducted in our lab on these mechanical concerns, it remains unpublished, making this thesis a critical step in formalizing and extending that foundational effort.

the safety
the 5 gauss line
the examples of safety
ASTM translational and torques studies
the current protocols to assess objects near MRI field
how this is different for us, size and suceptibility



ASTM F2052-21 offers a standardized technique for assessing magnetically induced displacement force on medical equipment in order to set safety standards for such situations. This standard states that a device is MR Conditional if the deflection force it encounters in the MRI environment is less than the force of gravity upon it. This is usually verified if the deflection angle from vertical is less than 45$^\circ$ \cite{stoianovici2024,astmF2052}.